% ===================================================================
%  TAULA N° 9 — La Montanha e l'Estacion termala. — Lo Bòsc e la Caça.
%  Traduction 2026 d'apres texto/io, controlee sur texto/fr et
%  texto/ca. Consigne : voir texto/oc/10-taula-01.tex.
% ===================================================================

%%K t09-apar-1 apar
\VUsaut{4.0mm}
\VUtitre{15.0pt}{60}{TAULA N° 9}
\VUsaut{3.0mm}

%%K t09-tit-1 sub
\VUcentre{12.0pt}{0}{\textit{Primièra scèna.}}
\VUsaut{3.0mm}

%%K t09-tit-2 sub
\VUpk{11.4pt}{40}{La Montanha. --- La Vila termala.}
\VUsaut{3.0mm}

%%K t09-c1-01-1 p
1. --- Fa uèch jorns que soi a Pratz. Pòdi pas exprimir tota l'admiracion que
sentiguèri quand me trobèri davant la meravilhosa \VUgras{sèrra}
\textsuperscript{(1)} dels Pirenèus, que sa \VUgras{cresta} \textsuperscript{(2)}
se retalha sul fons blau del cèl coma un feston gigantesc.

%%K t09-c1-02-1 p
2. --- M'imaginavi pas que los \VUgras{cims} \textsuperscript{(3)} pirenencs
arribèsson a una \VUgras{nautor} tan granda, e demorèri estabosit davant los
magnifics \VUgras{glacièrs} \textsuperscript{(5)} e los \VUgras{pics}
\textsuperscript{(4)} cobèrts de nèu ont rodolan d'\VUgras{avalancas} tan
terriblas.

%%K t09-tit-3 sub
\VUsaut{3.0mm}
\VUpk{10.2pt}{40}{Païsatge de montanha.}
\VUsaut{3.0mm}

%%K t09-c1-03-1 p
3. --- L'endeman meteis de mon arribada anèri al vièlh \VUgras{volcan}
\textsuperscript{(6)} atudat qu'es contra Pratz. Suls bòrds arids e nusses del
\VUgras{cratèr} se veson encara fòrça plan las traças del passatge de la
\VUgras{lava} que coriá, i a mantun sègle, pel \VUgras{pendís de la montanha}
\textsuperscript{(7)}. Capitèri de trobar, sus un dels \VUgras{versants},
qualques tròces d'aquela lava.

%%K t09-c1-04-1 p
4. --- Pratz es a la frontièra, e la \VUgras{fortalesa} \textsuperscript{(8)} que
defend lo \VUgras{pòrt} \textsuperscript{(27)} que separa França d'Espanha
remembra tot a fach aqueles vièlhs castèls de las ribas del Rin que semblan
espinchar una \VUgras{presa} de naut de lor ròca.

%%K t09-c1-05-1 p
5. --- Una \VUgras{agla} \textsuperscript{(9)} enòrma, qu'a son nis pas luènh del
\VUgras{fòrt} \textsuperscript{(8)}, atira sovent mon atencion. Aqueste
\VUgras{aucèl de rapina}, de \VUgras{bèc} poderós, pòrta sovent a sos pichons un
anhèl o un cabrit tengut dins sas \VUgras{arpas}. Sas \VUgras{alas} son tan
grandas que pòt planar longtemps amb sa carga pesuca.

%%K t09-tit-4 sub
\VUsaut{3.0mm}
\VUpk{10.2pt}{40}{L'Excursionista.}
\VUsaut{3.0mm}

%%K t09-c1-06-1 p
6. --- Lo passejament mai agradiu d'aicí es l'excursion a la \VUgras{cascada}
\textsuperscript{(10)} de Cau. Lo \VUgras{saut d'aiga} tomba en revolums e puèi se
pèrd dins un polit \VUgras{riu} dins la \VUgras{sapinèda} \textsuperscript{(11)}
situada a l'oèst de la vila. Pugèrem per aquel bòsc fins a l'\VUgras{observatòri
meteorologic} \textsuperscript{(12)} e, de naut del \VUgras{miralhador},
abracèrem lo \VUgras{panorama} entièr del país. Lo \VUgras{païsatge} es superbe e
la \VUgras{natura} sembla aver prodigat a aqueste país totes los tresaurs que ne
dispausa.

%%K t09-c1-07-1 p
7. --- Per davalar prenguèrem lo \VUgras{funicular} \textsuperscript{(13)} e nos
arrestèrem dins una pichona \VUgras{estacion} contra las \VUgras{roïnas}
\textsuperscript{(14)} d'un vièlh \VUgras{castèl feudal} abitat antan pels
senhers de Pratz.

%%K t09-c1-08-1 p
8. --- Paure castèl! Qué deu pensar en veire a sos pès la coqueta \VUgras{vila
termala} \textsuperscript{(15)}, qu'uèi es una \VUgras{estacion termala} a la
mòda?

%%K t09-c1-08-2 p
Lo \VUgras{clima} es tan agradiu e l'\VUgras{aiga minerala} tan benefica que la
\VUgras{cura} que se seguís al \VUgras{sanatòri} \textsuperscript{(17)} es un
veritable plaser.

%%K t09-tit-5 sub
\VUsaut{3.0mm}
\VUpk{10.2pt}{40}{Una agradiva vila termala.}
\VUsaut{3.0mm}

%%K t09-c1-09-1 p
9. --- Autrament, s'es fach tot per que l'estada siá agradiva. Se bastiguèt un
grand \VUgras{viaducte} \textsuperscript{(16)} per i poder tirar una via ferrada;
lo \VUgras{pargue} \textsuperscript{(18)} del \VUgras{establiment termal}, ont se
donan de \VUgras{concèrts}, es arrengat d'un biais encantador. Se son bastidas
mantunas graciosas \VUgras{vilas} \textsuperscript{(19)} a l'entorn del
\VUgras{casinò} \textsuperscript{(21)}, e una \VUgras{caussada}
\textsuperscript{(22)} fòrça plan suenhada fa aisidissimas las comunicacions.

%%K t09-c1-10-1 p
10. --- Un simple \VUgras{talús} \textsuperscript{(23)} separa la \VUgras{rota}
\textsuperscript{(20)} de la \VUgras{val} \textsuperscript{(25)} pròpriament
dicha, al fons de la quala s'espandís un polit \VUgras{lac}
\textsuperscript{(26)} qu'alimenta un \VUgras{riu} \textsuperscript{(24)} fòrça
clar.

%%K t09-c1-11-1 p
11. --- De l'autre costat de la val s'esplecha una \VUgras{mina}
\textsuperscript{(28)} de fèrre. Anèri mantun còp veire davalar los
\VUgras{minaires} pel \VUgras{potz}, e dintrèri quitament dins la
\VUgras{galariá} ont passan lo jorn a traire lo \VUgras{minerau} que conten lo
\VUgras{metal}.

%%K t09-c1-12-1 p
12. --- Se bastiguèt una importanta \VUgras{fonderiá} contra la mina per profitar
sulcòp de las riquesas que se'n tiran. Lo \VUgras{naut forn}
\textsuperscript{(29)}, escalfat amb lo \VUgras{carbon de pèira} qu'abonda aicí,
permet d'obténer de \VUgras{fonta} d'un biais rapid e a bon mercat.

%%K t09-c1-13-1 p
13. --- Pas luènh de la mina se dobrís una \VUgras{peirièira}
\textsuperscript{(30)}. Es ni \VUgras{granit}, ni \VUgras{pòrfir}, ni quitament
\VUgras{gres} çò que lo vièlh \VUgras{peirièr} arranca amb son \VUgras{pic}, mas
de simpla \VUgras{pèira de talha} per bastir d'ostals.

%%K t09-c1-14-1 p
14. --- Un dels mai bèles ornaments d'aqueste país es lo \VUgras{sap}
\textsuperscript{(31)}. Pertot --- al fons d'un \VUgras{ravin}
\textsuperscript{(32)}, a la riba d'un \VUgras{gorg} \textsuperscript{(33)},
contra un \VUgras{avenc} \textsuperscript{(34)} o un precipici ---, sas finas
agulhas i pausan lor nòta verda.

%%K t09-tit-6 sub
\VUsaut{3.0mm}
\VUpk{10.2pt}{40}{Caminar a pè.}
\VUsaut{3.0mm}

%%K t09-c1-15-1 p
15. --- Profiti mas vacanças per far de longs passejaments. Los \VUgras{camins}
\textsuperscript{(35)} que serpentejan per la montanha permeton de percórrer, sens
s'avisar de la distància, mantun \VUgras{quilomètre} e sovent mantuna
\VUgras{lèga}.

%%K t09-c1-15-2 p
De \VUgras{bòrnas} \textsuperscript{(36)} e d'\VUgras{indicadors}
\textsuperscript{(37)} jalonan los camins, ont òm rescontra sovent un jove
\VUgras{montanhòl} \textsuperscript{(38)} amb lo \VUgras{sac}
\textsuperscript{(40)} al costat, que pòrta una \VUgras{marmòta}
\textsuperscript{(39)}, o quauque \VUgras{gitano} \textsuperscript{(41)} que sa
\VUgras{bonèta de pèl} \textsuperscript{(42)} contrasta d'un biais singular amb
son vièlh \VUgras{capuchon} \textsuperscript{(43)} esquiçat. Mena a la
\VUgras{cadena} un \VUgras{ors} \textsuperscript{(44)} domdat.

%%K t09-tit-7 sub
\VUpk{10.2pt}{40}{Escalar de montanhas.}
\VUsaut{3.0mm}

%%K t09-c1-16-1 p
16. --- Gaireben cada jorn vau amb un \VUgras{guida} \textsuperscript{(48)} far una
\VUgras{ascension}, e soi fòrça fièr quand, amb lo \VUgras{sac de montanha}
\textsuperscript{(47)} a l'esquina e lo \VUgras{baston alpin} a la man, coma un
veritable \VUgras{torista} \textsuperscript{(46)}, escali las \VUgras{ròcas}
\textsuperscript{(45)}.

%%K t09-c1-17-1 p
17. --- De naut d'una montanha, los abitants de la plana semblan pas mai grands
que de formigas; un \VUgras{tropèl de fedas} \textsuperscript{(50)} que pastura
dins un \VUgras{pastenc} \textsuperscript{(49)} sembla una joguina d'enfant. Un
\VUgras{pin} \textsuperscript{(51)}, o quitament una \VUgras{cabana de fusta}
\textsuperscript{(52)}, son tot bèl just una taca dins la verdura.

%%K t09-c1-18-1 p
18. --- Pendent aquelas excursions nos crosam sovent sul \VUgras{sendarèl}
\textsuperscript{(53)} amb un \VUgras{caçaire d'isards} \textsuperscript{(54)} que
pòrta a l'espatla l'animal que ven de tuar.

%%K t09-c1-19-1 p
19. --- Ai pausat dins mon erbièr una branca de \VUgras{falguièra}
\textsuperscript{(55)} fòrça curiosa e un polit ramelet ròse de \VUgras{bruga}
\textsuperscript{(57)}, qu'ai culhit a l'intrada d'una pichona \VUgras{bauma}
\textsuperscript{(56)} dins mon darrièr passejament.

%%K t09-tit-8 sub
\VUsaut{3.0mm}
\VUcentre{12.0pt}{0}{\textit{Segonda scèna.}}
\VUsaut{3.0mm}

%%K t09-tit-9 sub
\VUpk{11.4pt}{40}{Lo Bòsc. --- La Caça.}
\VUsaut{3.0mm}

%%K t09-c2-01-1 p
1. --- Ièr passèri un jorn deliciós dins lo bòsc. L'activitat que règna entre los
animals e los \VUgras{insèctes} \textsuperscript{(5)} que lo poblan
m'interessèt vivament.

%%K t09-c2-01-2 p
L'\VUgras{aranha} \textsuperscript{(1)} que teis sa \VUgras{tela}
\textsuperscript{(2)} entre doas brancas per agafar la \VUgras{mosca}
\textsuperscript{(3)} que passa, la \VUgras{cagaraula} \textsuperscript{(4)}
que pòrta penosament sa closca, lo cèrvi volant que cèrca sa presa, la formiga
diligenta que s'afana contra lo \VUgras{formiguièr} \textsuperscript{(6)}: tot
aquel pichon mond anava, veniá e trabalhava amb una animacion extraordinària.

%%K t09-c2-02-1 p
2. --- Una \VUgras{sèrp} \textsuperscript{(7)}, qu'al començament prenguèri per
una \VUgras{vipèra} mas qu'èra pas qu'una simpla \VUgras{colòbra}, èra amagada
jos un tronc d'arbre e de temps en temps sortissiá son caponet prim, que semblava
totjorn prèst a mòrdre. Davant ieu, una gròssa \VUgras{limaça}
\textsuperscript{(8)} rampava pesucament aprèp aver devorat una de las saborosas
\VUgras{majofas} \textsuperscript{(11)} que i creisson en abondància.

%%K t09-c2-03-1 p
3. --- Lo sòl èra tapissat de \VUgras{flors del bòsc}: de \VUgras{primavèras}
\textsuperscript{(9)}, de \VUgras{pervencas} \textsuperscript{(10)} e fòrça
autras plantas que coneissiái pas florissián entre las \VUgras{tòfas d'èrba}
\textsuperscript{(12)}.

%%K t09-c2-04-1 p
4. --- Dins aqueste país, la \VUgras{trufa} es gaireben tan comuna coma lo
\VUgras{bolet} \textsuperscript{(14)}, e un \VUgras{porcelet}
\textsuperscript{(13)} d'un garda forestièr cercava amb un aire golut se ne
trobava qualqu'una.

%%K t09-tit-10 sub
\VUsaut{3.0mm}
\VUpk{10.2pt}{40}{La caça furtiva.}
\VUsaut{3.0mm}

%%K t09-c2-05-1 p
5. --- Assistiguèri a la \VUgras{fin} d'una scèna que me divertiguèt fòrça.
Gràcia a l'odorat de son \VUgras{gos basset} \textsuperscript{(15)}, lo
\VUgras{garda forestièr} \textsuperscript{(16)} venia de susprene un
\VUgras{caçaire furtiu} \textsuperscript{(22)} al moment ont aqueste s'aprestava a
emportar la \VUgras{caça} \textsuperscript{(17)} qu'aviá tuada. Aimant mai
abandonar sa presa que se daissar arrestar, lo furtiu aviá fugit, e
\VUgras{lèbre} \textsuperscript{(18)}, \VUgras{cèrva} \textsuperscript{(19)},
\VUgras{cabiròl} \textsuperscript{(20)} e \VUgras{singlar}
\textsuperscript{(21)} demoravan sus l'èrba, per la jòia del garda.

%%K t09-c2-06-1 p
6. --- La varietat d'arbres que sarra lo bòsc es espantanta. Los \VUgras{fages}
\textsuperscript{(23)} e los \VUgras{beçes} \textsuperscript{(26)}, los
\VUgras{pins}, que donan la \VUgras{rosina} \textsuperscript{(27)}, los
\VUgras{garrics} \textsuperscript{(29)} e fòrça mai entremesclan lors brancas.

%%K t09-c2-06-2 p
L'\VUgras{èdra} \textsuperscript{(24)}, que s'agafa al tronc dels arbres auts,
dona a la \VUgras{sèlva} \textsuperscript{(25)} un verd brilhant que s'unís
agradivament al ton mai clar e mai doç de la \VUgras{mossa}
\textsuperscript{(31)} ont lo \VUgras{pigòt verd} \textsuperscript{(30)} cèrca
d'insèctes.

%%K t09-tit-11 sub
\VUsaut{3.0mm}
\VUpk{10.2pt}{40}{Païsatge de bòsc.}
\VUsaut{3.0mm}

%%K t09-c2-07-1 p
7. --- Los vièlhs garrics m'encantèron. Lors gròssas \VUgras{raiças}
\textsuperscript{(32)} tòrtas, mièg sortidas de la tèrra, lor donan una esplendida
aparéncia de fòrça e de vigor, e me demandi cossí lo \VUgras{proprietari}
\textsuperscript{(35)} se pòt resignar a los far abatre e a los liurar a la
\VUgras{rèssa} \textsuperscript{(34)} del \VUgras{bosquetièr}
\textsuperscript{(33)}. Los paures \VUgras{troncs} \textsuperscript{(36)},
despolhats de sa \VUgras{rusca} \textsuperscript{(37)} o fenduts per la
\VUgras{destral} \textsuperscript{(38)} del \VUgras{jornalièr}
\textsuperscript{(39)}, me fan pietat.

%%K t09-c2-08-1 p
8. --- A prèp d'aquí, un vièlh \VUgras{carbonièr} \textsuperscript{(41)} aviá
preparat amb sa \VUgras{pala} un \VUgras{molon} \textsuperscript{(42)} de carbon,
e una vièlha emportava un \VUgras{fais} \textsuperscript{(40)} de \VUgras{lenha
mòrta} fach amb los ramèls dels arbres abatuts.

%%K t09-tit-12 sub
\VUsaut{3.0mm}
\VUpk{10.2pt}{40}{La Caçada.}
\VUsaut{3.0mm}

%%K t09-c2-09-1 p
9. --- Voliái anar me pausar una estona a l'\VUgras{ostal del garda forestièr}
\textsuperscript{(46)}, mas quand arribèri contra l'\VUgras{estanh}
\textsuperscript{(43)}, ont nada un bèl \VUgras{cinne} \textsuperscript{(45)},
m'avisèri qu'una \VUgras{inondacion} \textsuperscript{(44)} recenta aviá cambiat
lo camin en un veritable \VUgras{palun}. Me calguèt far un torn e, en passar
contra lo \VUgras{cèdre} \textsuperscript{(49)}, los \VUgras{pibols}
\textsuperscript{(48)} e l'\VUgras{olme} \textsuperscript{(47)} que son a la
lisièra del bòsc, vegèri sortir d'una \VUgras{bartassa} \textsuperscript{(54)} un
magnific \VUgras{cèrvi} \textsuperscript{(50)} perseguit per una \VUgras{mainada
de gosses} \textsuperscript{(51)} enrabiada, qu'excitava en mai lo \VUgras{còrn}
\textsuperscript{(53)} dels \VUgras{caçaires a caval} \textsuperscript{(52)}. Lo
paure animal èra espotit. Venguèt tombar moribond a qualques passes de ieu.

%%K t09-c2-10-1 p
10. --- Lo filh del garda forestièr, atirat pel bruch de la \VUgras{caçada},
sortiguèt de l'ostal e tornèrem totes dos al bòsc. Aviá vista una \VUgras{agaça}
\textsuperscript{(56)} sus la branca d'un vièlh garric. Pensava que son nis deviá
èsser amagat dins lo \VUgras{fulhatge} \textsuperscript{(58)} e lo voliá prene.

%%K t09-c2-10-2 p
Agil coma un \VUgras{esquiròl} \textsuperscript{(61)}, s'enfilèt de
\VUgras{branca} \textsuperscript{(55)} en branca, mas trobèt res e capitèt pas que
de far volar una vièlha \VUgras{grala} \textsuperscript{(60)} pausada sus un
\VUgras{ramèl gròs} \textsuperscript{(57)}.
\VUsaut{4.0mm}
\VUfilet{22mm}
